\begin{claim}{AC.37}{Proved: a free order-117 group and its precise limitation}
On the Fermat curve \(D/\mathbf F_{4096}\) of AC.36,
choose a primitive thirty-ninth root \(\xi\) and put
\[
 s(X:Y:Z)=(X:\xi Y:\xi^{17}Z),
 \qquad \tau(X:Y:Z)=(Y:Z:X).
\]
The constant nonabelian group
\[
                 G'=\langle s,\tau\rangle
                    \simeq C_{39}\rtimes_{16}C_3
\]
has order \(117\) and acts freely. Its quotient \(C_7=D/G'\)
has genus seven and zeta numerator
\[
 P_{C_7}(T)=(1-64T)^2(1+64T)^4
                         (1-115T+4096T^2)^4.
\]
Thus \(J_{C_7}(\mathbf F_{4096})[2^\infty]
\simeq(\mathbf Z/2)^4\), and \(D/C_7\) again has relative
class-number valuation twenty.
The bounded-exponent criterion for this cover applies
after splitting the pointed fibre over every extension
of \(\mathbf F_{4096}\) of odd or twice-odd degree:
the relevant exponents are respectively two and eight,
and \(117>(8-1)^2\).

There is also a more direct (R) witness for this particular
quotient. The subgroup \(\langle s^3,\tau\rangle\) is
normal of index three and is the order-39 group of AC.36,
after choosing its generator of order thirteen.
Writing \(C_{19}=D/\langle s^3,\tau\rangle\), the cover
\(C_{19}\to C_7\) is etale cyclic of degree three, with
relative zeta numerator and relative class number
\[
                  (1+64T)^{24},\qquad 65^{24}.
\]
Its new two-rank is zero, so (R) persists over every
extension, and every prescribed base point becomes
split after an additional extension of degree at most
three. Accordingly, the smaller genus of this quotient
does not enlarge the set of curves reached by (R).
\end{claim}
\begin{proof}
Projectively,
\[
       \tau s\tau^{-1}=s^{16},\qquad
       16^3\equiv1\pmod{39},\qquad
       1+16+16^2\equiv0\pmod{39}.
\]
The diagonal generator has order \(39\), and no nonidentity
power of \(\tau\) is diagonal, giving the stated order
and semidirect product. The pairwise differences of the
three diagonal exponents \(0,1,17\) are \(1,16,17\),
all units modulo \(39\). Hence a nonidentity diagonal
element has distinct eigenvalues and fixes only the
coordinate vertices.

For an element \(s^a\tau^b\), with \(b=1\) or \(2\),
use its displayed monomial matrix. Its cube is
\(\xi^{18a}I\). If \(\lambda\) is an eigenvalue,
\[
              \lambda^3=\xi^{18a},\qquad
              \lambda^{39}=\xi^{234a}=1.
\]
All eigenvector coordinates are nonzero; their ratios
are products of powers of \(\lambda\) and \(\xi\),
so their thirty-ninth powers are equal. The Fermat
expression at an eigenpoint is therefore a nonzero
scalar times \(1+1+1=1\). There are no fixed points.
Riemann--Hurwitz gives
\(g(C_7)=1+(703-1)/117=7\).

For the Jacobi character indexed by \((a,b,c)\),
\(c=-a-b\), the diagonal invariance condition is
\[
                   b+17c\equiv22a+23b\equiv0\pmod{39}.
\]
The augmented exact program finds \(38\) characters:
the same two cycle-fixed characters as in AC.36 and
twelve cycles of length three. The quotient eigenvalues
are \(64\) twice, \(-64\) four times, and the two
roots of \(X^2-115X+4096\) four times each.
This proves the displayed numerator. The actual
integral exponent follows from the scalar unit-root
argument of AC.36.

Every element of \(G'\) has order in
\(\{1,3,13,39\}\). A rational fibre therefore splits
after an odd extension of one of those degrees.
Such an extension does not change the two-adic valuation
of the original constant-field degree. Hence the exact
exponent remains two for odd degrees and eight for
twice-odd degrees, as stated.

Finally, the map \(G'\to C_3\) sending \(s\) to a
generator and \(\tau\) to zero is well-defined because
\(16\equiv1\pmod3\). Its kernel is
\(\langle s^3,\tau\rangle\); the order-thirteen
generator \(s^6\), after multiplying coordinates by
a common scalar, is the diagonal generator in AC.36.
Dividing the two explicitly proved zeta numerators gives
\((1+64T)^{24}\). Its roots all have positive two-adic
valuation; alternatively its value over every extension
is an odd power of \(1-(-64)^m\). Thus (R) persists
under the extension used to split any given fibre.
\end{proof}
